% s01 Introduction and main results
\section{Introduction and main results}

\subsection{King's construction}

King \cite{K1,K2} constructed the three-dimensional $U(1)$ Higgs model, scalar electrodynamics with a
massive photon, as a continuum and infinite-volume limit of lattice models. The lattice theory lives on
$\Lk=\eps_k\Z^3$ with spacings $\eps_k=5^{-k}$; its action has a photon $A$ in Feynman form with mass
$\lmu>0$, an $N$-component scalar $\phi$ with quartic coupling $\lambda>0$ and bare mass $m>0$, minimally
coupled with charge $e>0$, and one mass counterterm $\dm_{\eps_k}$. The counterterm is King's free-boundary
one: it is computed with infinite-lattice propagators, so it is one number per spacing, the same at every
site of every box and every torus. The ultraviolet problem is settled in \cite{K1} with Ba{\l}aban's
renormalisation group \cite{Ba1,Ba2,Ba3,Ba4} at block factor $5$: on a finite volume the effective
actions converge as $k\to\infty$. The infinite-volume problem is treated in \cite{K2} with the correlation
inequalities of Brydges, Fr\"ohlich and Seiler \cite{BFS1}. King's finite-volume states are
Neumann--Dirichlet states on the cubes $C_J=[-MJ,MJ)^3$, $M\ge1$ his block integer: the photon gradient
terms are kept only inside the cube (Neumann), and the scalar field vanishes on a boundary layer
(Dirichlet). The limit is iterated, first $k\to\infty$ at fixed $J$, then $J\to\infty$; the second limit
exists because the cube states are monotone in $J$ on a cone of observables, by \cite{BFS1}.

This paper proves the Osterwalder--Schrader axioms for King's printed state at fixed photon mass, with
full Euclidean invariance, and the corresponding statements for the joint field of the curvature and the
renormalised Wick square on a tube of small Wick amplitudes. It is self-contained in the following sense.
The existence of King's cube limits, the existence clause of his Theorem~2.1 \cite{K2}, is proved in the
existence part of the paper (\cref{sec:k01}, and \cref{sec:k02} for the
version with scalar sources): King's route, the comparison of two spacings at a common block scale, is run
in a ratio form in which the vacuum energies cancel, so that the convergence of the vacuum energies, which
his printed argument uses and which is not available on the cube, is not needed. The passage from the cube
limits to the axioms in \cite{K2} is asserted at three places, and at each we proceed differently. First,
King's Neumann--Dirichlet cube form is not reflection symmetric about any plane, and its natural
symmetrisation is not reflection positive (an exact counterexample at $e=0$, \Cref{rem:s07:R1}), so his
reflection positivity goes through a modified box state; that this box state is reflection positive is his
Proposition~3.3, which he derives from \cite[Cor.~5.4]{BFS1}, whose hypotheses do not cover the finite-box
Gaussian of the modified action. Proposition~3.3 is true, and it is proved here (\Cref{lem:s07:RPbox})
rather than derived from \cite[Cor.~5.4]{BFS1}. Second, the inequality by which he passes from the box
states back to the cube limit holds in the reverse direction for the class of observables his argument
needs: weakly, at every coupling, for the box state with the site-independent counterterm, and strictly at
$e=0$ (\Cref{lem:s07:R}); it is replaced by a two-sided sandwich between cube states with a margin of order
$\lmu^{-1}$ (\Cref{lem:s07:L,lem:s07:U}). Third, the continuum subsequences through which he reaches the
limit theory are built here (\Cref{rem:s12:K9}). With scalar sources, finally, his Theorem~2.1 needs a
vacuum counterterm that his Wick ordering does not contain; we work with a scheme $a$ and prove the
existence of the cube limits with sources in that scheme.

The authors' earlier paper \cite{SymPhase} constructs torus limits of the same model at block factor $2$
and proves the Osterwalder--Schrader axioms with cubic symmetry for them; its abstract measure theory (the
reconstruction of the measure, its ergodic decomposition and the weak clustering of the components) is used
here as it stands, cited by label. Rotation invariance is proved in the rotation part of this paper, whose
main theorem (\Cref{rot:thm:main}) precedes the main theorem of the paper: King's cubes are compared with periodic tori carrying
a smooth mass wall, the walled tori on two lattices whose orientations differ by $\arctan(3/4)$ are
compared, and King's state is identified as the lower envelope of the walled tori through a heavy-shell
decoupling estimate on King's block-spin flow with the shell mass inside the matter covariance.

\subsection{The two cube conventions}

King's cube state $\NDK_{k,J}$ (\Cref{def:s02:Kcube}) has its photon variables on $X_J=C_J\cap\Lk$ and its
scalar alive on $\{x\in X_J:\ x_\nu<MJ-\eps_k\ \text{for all }\nu\}$, that is, killed on the three top
face layers $T_J$ of the cube. The comparison cube state $\NDa_{k,J}$ (\Cref{def:s02:acube}) has the
scalar alive on all of $X_J$. The Wick-tube lemmas and the lattice comparisons of the rotation part are
written for the comparison convention, while King's printed state is on his own. The two are interlaced:
for observables of the cone class, with supports at $\ell^\infty$-distance $>1$ from the complement of the
cube,
\[
\avK{\Theta}\ \le\ \ava{\Theta}\ \le\ \langle\Theta\rangle^{K}_{k,J+1},
\]
and reversely for the opposite class (\Cref{lem:s04:G}). Hence both families have the same $J$-limit at
fixed $k$; the comparison family's fixed-cube limits inferior and superior are squeezed to King's printed
state $\ZKinf$ as $J\to\infty$, and every continuum-subsequence limit of the comparison family equals
$\ZKinf$ (\Cref{cor:s04:G2}). What does not transfer is the existence of the fixed-$J$ limit $k\to\infty$
itself: on one convention it does not imply it on the other. This is why the existence theorems are proved
on King's convention, where the printed state lives and where King's block-spin flow runs; the Wick tube,
proved on the comparison convention, is carried to King's convention by a transfer theorem
(\Cref{thm:s13:TWK}).

\subsection{The main theorem}

Sources are real $2$-forms $f_j\in C^\infty_c(\R^3;\Lambda^2)$ with coefficients $z_j$ and real scalar test
functions $h_j\in C^\infty_c(\R^3)$ with coefficients $w_j$; $F(f)$ is the curvature paired with $f$,
$W^a(h)$ the renormalised Wick square in the scheme $a$, and $g(z)=\sum_j(\Re z_j)f_j$,
$k(\zeta)=\sum_j(\Re w_j)h_j=k_+-k_-$, with the norms $\norm{\cdot}_2$, $\fournorm{\cdot}$,
$\wnorm{\cdot}$ and the tube $\Dtube$ of \Cref{def:s02:sources,def:s02:tube}. The Schwinger functions
$S_n$, $S_{n,p}$ are the Taylor coefficients of the limit functionals at $0$. A scheme $a$ is admissible
if $a=z$, the reference scheme, or $z_k-a_k$ converges (\Cref{def:s03:adm}). The following summary is
restated with every clause in full and proved as \Cref{thm:s13:main}.

\begin{theorem}[Main Theorem, summary]\label{thm:s01:main}
Let $e>0$, $\lambda,m,\lmu>0$ and $M\ge1$ be arbitrary.

\medskip\noindent\textbf{(A) The curvature field.} King's printed state
$\ZKinf(f,0)=\lim_{J}\lim_{k}\avK{e^{F(f)}}$ exists for every real $2$-form $f$, extends to an entire
function of the coefficients of every finite family, and satisfies:
\begin{enumerate}[label=\textup{(\roman*)},leftmargin=*]
\item \textup{(\OSz)} $\abs{\ZKinf(\sum_jz_jf_j,0)}\le e^{\norm{g(z)}_2^2}$ for all $z\in\C^N$, with no
dependence on $\lmu$, $e$, $\lambda$, $m$ or $M$, and
$\abs{\ZKinf(\sum_jz_jf_j,0)}\le\exp\big(\tfrac12\lmu^{-2}\norm{\sum_j(\Re z_j)\dstar f_j}_2^2\big)$;
for real test functions $\abs{S_n(f_1,\dots,f_n)}\le(\sqrt2e)^n(n!)^{1/2}\prod_j\norm{f_j}_2
\le(e\pi/2)^n(n!)^{1/2}\prod_j\fournorm{f_j}$; the $S_n$ are tempered distributions and satisfy the
regularity axiom \Ezz\ with $s=4$, the constant $\sigma_n=(e\pi)^n(n!)^{1/2}$ for complex
$f_j\in\cS$, and every $\beta>\tfrac12$.
\item \textup{(\OSo)} Invariance under the full Euclidean group $\Euc$: under all translations of
$\R^3$, under the cubic group $\Oh$ about every point, and under all rotations, $F$ transforming as a
$2$-form; the $S_n$ are Euclidean covariant.
\item \textup{(\OSt)} Reflection positivity about every plane $\{x\cdot n=c\}$, $n$ a unit vector,
$c\in\R$, in matrix form for exponentials with arbitrary complex curvature coefficients and in polynomial
form.
\item \textup{(\OSth)} $S_n$ is real and symmetric, and $S_n=0$ for odd $n$.
\item \textup{(The measure)} For every even $r\ge14$ there is exactly one Borel probability measure
$\mu_F$ on $\cH^{-r}(\R^3;\Lambda^2)$ with moments $S_n$ and characteristic functional
$\ZKinf(\mathrm i\,\cdot,0)$ on $\cS$; its image on $H^{-r}_{\rm loc}$ is determined by the characteristic
functional and by the moments. The laws of the localised lattice fields $\chi F_{\Lk}$ under $\NDK_{k,J}$
converge weakly on $H^{-r}_{\rm per}(Q')$, $Q'$ a box containing $\supp\chi$, along the full sequence
$k\to\infty$ at every fixed $J\ge J_0(\chi)$, and then as $J\to\infty$, to the law of $\chi F$ under
$\mu_F$.
\item \textup{(\OSf\ on the ergodic components)} $\mu_F$ is $\Euc$-invariant and $\mu_F=\int\mu_\xi\,d\nu(\xi)$
with translation-ergodic components; $\nu$-almost every component is $\Euc$-invariant, reflection positive
about every plane, clusters in the mean and under diagonal separation, has a unique joint vacuum in every
direction, satisfies \OSf\ in all its forms for every choice of the time direction, and has real symmetric
moments with a factorial growth bound with component constants; the decomposition is $\Oh$-invariant, each
component being a mixture of $k\in\{1,2,3,6\}$ ergodic measures. The vanishing of the odd moments is not
claimed for a component.
\end{enumerate}

\medskip\noindent\textbf{(B) The joint field on the Wick tube.} Let the couplings satisfy the sign clause
\Tthree\ of \Cref{hyp:s03:T3}, $\sigma_1(M_1)>0$ for some $M_1\in(0,m)$; in particular, let them lie in
one of the explicit regions of \Cref{prop:s14:T3}. Let $a$ be an admissible scheme. Then for every
family of sources supported in $Q_R=[-R,R]^3$ the joint functional of $(F,W^a)$ has a limit $\Ginf$ on
the tube $\Dtube$, with
\[
\sigma_+=s_1,\qquad H=\tfrac{s_1}{16},\qquad \kappa=\tfrac{C_U}{s_1},
\]
and $K_-$ determined by the cell constant $B_{\rm cell}$ (\Cref{eq:s11:constants}), where
$s_1\in(0,\sigma_1(M_1))$ is the amplitude of the extensive bound and $C_U$ its constant
(\Cref{eq:s10:I2,prop:s10:U}); the constants depend on $(e,\lambda,m,\lmu,M,a)$ and on the auxiliary choices $M_1$,
$\chi_0$, $s_1$, and on nothing else, and they are not uniform in $\lmu$. The joint Schwinger functions
$S_{n,p}$ satisfy \textup{(i)--(vi)} in the following form:
\begin{enumerate}[label=\textup{(\roman*)},leftmargin=*]
\item $\abs{\Ginf(\zeta)}\le\exp\big(\norm{g}_2^2+\kappa\norm{k_+}_{L^1}+K_-\wnorm{k_-}\big)$ on
$\Dtube$, and for real test functions
$\abs{S_{n,p}}\le e^{K'\sigma_0}(\sqrt2e)^n(n!)^{1/2}(e/\sigma_0)^p\,p!\prod_i\norm{f_i}_2\prod_j\wnorm{h_j}$
with $\sigma_0=H/4$ and $K'=\kappa\pi^2+K_-$ (a factor $2^{n+p}$ for complex test functions); \Ezz\ holds
with $s=4$ and every $\beta>1$.
\item Translations and $\Oh$ about every point hold as identities on the tube, $W^a$ transforming as a
scalar; there is no Euclidean clause.
\item Reflection positivity about every coordinate plane for families whose Wick amplitudes satisfy
$\amp(\zeta_i)<\sigma_W=\sigma_+/8$, in particular for all purely imaginary Wick coefficients, and the
polynomial form.
\item $S_{n,p}=0$ for odd $n$.
\item[\textup{(v), (vi)}] As in \textup{(A)} for the measure on $\cH^{-r}(\R^3;\Lambda^2\oplus\R)$,
$r\ge14$ even, with the scheme-shifted lattice field in the convergence clause, and with the weak form of
\OSf\ (ergodicity under $\R^3$ and clustering in coordinate directions) on $\nu$-almost every component;
there is no Euclidean clause.
\end{enumerate}
The same statements hold for the tube object $G^K_\infty$ of King's printed state on his own convention
(\Cref{thm:s13:TWK,prop:s13:WK}).
\end{theorem}

The $\lmu$-free bound in (A)(i) is absolute: it depends on none of the parameters of the model. The
$\lmu^{-2}$ bound is the finite-volume Gaussian bound carried to the limit. Part (B) is the headline form
of the Wick leg; the extension of (B) to arbitrary Wick amplitudes, under the hypothesis \Wtwo\ at every
amplitude, with the only Euclidean clause of the Wick leg, is stated separately as \Cref{thm:s23:FormA}
among the open problems.

\subsection{The massless limit}

The photon mass $\lmu$ is fixed throughout: King's construction, and every statement above, is at
$\lmu>0$, and the constants of the Wick leg are not uniform in $\lmu$. The limit of vanishing photon mass,
along a family whose counterterm and scheme are frozen at a reference mass, is not treated in this paper;
it is listed among the open problems as later work.

\subsection{The coupling range}

The curvature leg holds on all of King's range: $e>0$ and $\lambda,m,\lmu>0$ arbitrary. Every smallness
condition in \cite{K1,K2,Ba1,Ba2,Ba3,Ba4} is a cap on the running couplings at a stopping scale that is
chosen after the point $(e,\lambda,m,\lmu)$ and the datum, so the quantifier over the couplings is the one
over the sources, and no sentence ``$e,\lambda$ sufficiently small'' occurs
(\Cref{lem:s14:CW,rem:s14:quantifier}). No relation $\lmu\le m$ is used or needed. The only coupling
condition in the paper is the one-sided sign clause \Tthree\ of the Wick leg, $\sigma_1(M_1)>0$ for some
$M_1\in(0,m)$; it fails exactly when $m$ minimises King's free-boundary bare coefficient $\alpha_\infty$ on
$(0,m]$, and to one-loop order it holds if and only if
\[
e^2<e_c^2(\lambda,m,\lmu)=\max\Big\{\frac{8\pi(m+\lmu)^2}{m+2\lmu}\Big(1+\frac{\lambda(N+2)}{2\pi m}\Big),\
4\pi(m+\lmu)\Big(1+\frac{\lambda(N+2)}{\pi m}\Big)\Big\}
\]
(\Cref{eq:s14:ec2}), so that a half-line of couplings fails it at one loop. Beyond one loop, \Tthree\ is
proved on explicit coupling regions, at $\lmu=m$ on a region containing $0\le\lambda,e^2\le0.1m$ and at
every $\lmu\in(0,m]$ on explicit boxes (\Cref{prop:s14:T3}); the Wick leg is stated on those regions, and
more generally wherever \Tthree\ holds. The mass clause in Ba{\l}aban's upper bound, a cap on the
dimensionless scalar mass at the stopping scales where his Proposition~3.1 is used, is part of that
proposition's statement and is one more one-sided cap on the stopping scale, not a condition on the
couplings (\Cref{rem:s14:BALm}).

\subsection{What is proved where}

The inputs of the main theorem are collected in \Cref{sec:s03}: the existence theorems for King's cube
limits without and with scalar sources (\Cref{hyp:s03:EC,hyp:s03:W1}), proved in the existence part; the
rotation theorem (\Cref{hyp:s03:RC}), proved in the rotation part; Ba{\l}aban's torus stability bounds and
the $c$-number lemma in the form the Wick leg uses (\Cref{hyp:s03:T1,hyp:s03:T2}); and the coupling
condition \Tthree. Everything else is proved in the sections between. Open are only the optional
extensions: the Wick sector at arbitrary amplitudes (\Wtwo, Form~(A), \Cref{thm:s23:FormA}), the relation
of King's state to the torus states of \cite{SymPhase} and its uniqueness and ergodicity, the limit of
vanishing photon mass, and non-triviality; they are collected in the final section.

\subsection{Structure of the paper}

The model, King's cube, the comparison cube, the sources and the standing lemmas (the correlation
inequalities of \cite{BFS1}, with the two inequalities that are only sketched there proved in full, Gaussian
domination, cutting and the elementary facts) come first, then the inputs with their provenance. The
interlacing of the two cube conventions (\Cref{lem:s04:G,cor:s04:G2}) and its three geometry margins
follow, and then the existence part: King's existence theorem on the cube (\Cref{sec:k01}) and its version
with scalar sources. The curvature sector is then built up in the order of the axioms: regularity with the
absolute constants (\Cref{thm:s05:C,thm:s05:CK}), translations, the cubic group and the rotation clause
(\Cref{thm:s06:S1K,thm:s06:O3}), reflection positivity through the modified box and its passage to the
limit with \OSth\ and the polynomial form (\Cref{lem:s07:RPbox,thm:s08:OS2C,prop:s08:OS3,prop:s08:poly}),
and the measure with its ergodic decomposition (\Cref{thm:s09:S2K}). The Wick tube is treated in two
sections, the torus inputs and the extensive bound (\Cref{prop:s10:RP,prop:s10:U}) and the uniform tube
bound with its limits (\Cref{thm:s11:W,lem:s11:TWtube}); continuum subsequences and the states they
define, with no existence input, form a section of their own (\Cref{thm:s12:LIM}). The rotation part
follows: the geometry of the two lattices and the walled tori, the admissible schemes and the full object,
the torus comparison and its reduction to one-sided bounds at one spacing, the closing estimate and the
invariance theorems, and the heavy shell in two sections; a short section collects the corrections to
King's papers that it uses. The main theorem is assembled in \Cref{thm:s13:main}, and the coupling range
is established in \Cref{lem:s14:CW,prop:s14:T3}. The paper ends with the open problems.
